%---------------------------------------------------------
%	2. PROBLEM STATEMENT & OBJECTIVES
%---------------------------------------------------------
\section{Problem}

\begin{frame}
    \frametitle{Problem Statement}
    \framesubtitle{Night photography degrades images in three ways at once}

    \begin{columns}[c]
        \begin{column}{0.52\textwidth}
            \small
            Small phone sensors + fixed aperture $\Rightarrow$ \textbf{long exposures} at night:
            \begin{itemize}
                \item \textbf{Low light}: dark, low-contrast images
                \item \textbf{Noise}: shot / read noise of the sensor
                \item \textbf{Motion blur}: hand shake during the exposure
            \end{itemize}
            Image formation model:
            \[
                y = \gamma\big(x \otimes k\big) + n
            \]
            {\footnotesize $y$: captured image, $x$: clean scene, $k$: blur kernel, $n$: noise, $\gamma$: dynamic-range loss}

            \textbf{Goal:} recover $x$ from a single $y$ --- brighten, denoise \emph{and} deblur in one model.
        \end{column}
        \begin{column}{0.46\textwidth}
            \centering
            \includegraphics[height=0.31\textheight]{../../assets/teaser/0085_low.png}\\[-0.5mm]
            {\scriptsize Input: low light + blur + noise}\\[1.5mm]
            \includegraphics[height=0.31\textheight]{../../assets/teaser/0085_darkir.png}\\[-0.5mm]
            {\scriptsize NightClear output}
        \end{column}
    \end{columns}
\end{frame}

\begin{frame}
    \frametitle{Why Data Science? Objectives \& KPIs}

    \begin{columns}[t]
        \begin{column}{0.47\textwidth}
            \small
            \textbf{Why a data-driven model}
            \begin{itemize}
                \item The needed brightening is image-dependent: \textbf{$\times$2.8 to $\times$30} across LOLv2-real pairs (our EDA) --- no fixed curve fits all.
                \item Chaining two specialised models fails: \textbf{14.7--18.3\,dB} vs.\ \textbf{27.0\,dB} for one joint model.
                \item 10k+ paired examples (LOL-Blur) let a network learn the inverse mapping directly.
            \end{itemize}
        \end{column}
        \begin{column}{0.51\textwidth}
            \small
            \textbf{KPIs} {\footnotesize (target = beat LEDNet, the previous joint method)}\\[1mm]
            \footnotesize
            \begin{tabular}{lll}
                \toprule
                \textbf{Metric (LOLBlur)} & \textbf{Target} & \textbf{NightClear-m} \\
                \midrule
                PSNR $\uparrow$  & $> 25.74$ dB & \textbf{27.00} $\checkmark$ \\
                SSIM $\uparrow$  & $> 0.850$    & \textbf{0.883} $\checkmark$ \\
                LPIPS $\downarrow$ & $< 0.224$  & \textbf{0.162} $\checkmark$ \\
                Compute          & $< 10$ GMACs & \textbf{7.25} $\checkmark$ \\
                Model size       & $< 5$ M params & \textbf{3.31} $\checkmark$ \\
                Deployment       & usable app   & Gradio UI $\checkmark$ \\
                \bottomrule
            \end{tabular}
        \end{column}
    \end{columns}
\end{frame}
